\documentclass[12pt,a4paper,twoside,spanish]{report}
\usepackage{graphicx}
\usepackage{array}
\usepackage{enumerate}
\usepackage{graphicx, subfigure}
\usepackage[spanish,activeacute]{babel}
\usepackage[T1]{fontenc}
\usepackage{lmodern}
\usepackage[utf8]{inputenc}
\usepackage[left=2.5cm,top=3.5cm,right=2.5cm,bottom=2.5cm]{geometry}
\usepackage{fancyhdr}
\usepackage{setspace}
\usepackage{url}
\usepackage{sectsty}\allsectionsfont{\usefont{OT1}{phv}{bc}{n}\selectfont}
\setcounter{secnumdepth}{5}
\usepackage{ccaption}
   \captionnamefont{\bfseries\small}
   \captiontitlefont{\small}
   \captiondelim{\hspace{1em}}
   \hangcaption
\usepackage{multirow}
%\usepackage{longtable,ltcaption}
\usepackage{pdfpages} % to import PDF pages


\usepackage{relsize}
\usepackage{framed}
\usepackage{amsmath}
\usepackage{amsfonts}
\usepackage{color}
%\usepackage[ruled,linesnumbered]{algorithm2e}
\usepackage[ruled, noend]{algorithm2e}
\usepackage[table]{xcolor}% http://ctan.org/pkg/xcolor

\usepackage{listings}
\definecolor{backgroundcode}{rgb}{0.95,0.95,0.95}


%referencias de la pagina
\usepackage{hyperref}

%% propiedades del pdf
\hypersetup{
    pdfstartview={FitH},  
    pdftitle={Gu\'ia pr\'actica $4$ - $\the\year$},
    pdfauthor={David N. Campo - Daniel Ambort},     
    pdfsubject={Gu\'ia pr\'actica $4$ - Programaci\'on Competitiva \the\year  -- Complejidad},
    pdfcreator={David N. Campo},   
    pdfproducer={David N. Campo}, 
    pdfkeywords={algorithms, c++, acm icpc, UVa, SPOJ, programming}, 
    colorlinks,
    citecolor=black,
    filecolor=black,
    linkcolor=black,
    urlcolor=blue,
}

% Estilo de página %%%%%%%%%%%%%%%%%%%%
\fancyhead{}
\fancyhead[LO,RE]{\vfill{\footnotesize Programaci\'on Competitiva \the\year - Gu\'ia pr\'actica 4 -- Complejidad}}
\fancyhead[RO,LE]{\includegraphics[scale=0.15]{img/logo_UTN_color.pdf}}

\fancyfoot{}
\fancyfoot[RO]{\hfill\includegraphics[scale=0.9]{img/footlineR.pdf}\newline\footnotesize Página \textbf{\thepage}}
\fancyfoot[LE]{\includegraphics[scale=0.9]{img/footlineL.pdf}\newline\footnotesize Página \textbf{\thepage}}

\pagestyle{fancy}
\fancypagestyle{plain}{
   \fancyhead{}
	\fancyhead[LO,RE]{\vfill{\footnotesize Programaci\'on Competitiva \the\year - Gu\'ia pr\'actica 4 -- Complejidad}}
   \fancyhead[RO,LE]{\includegraphics[scale=0.15]{img/logo_UTN_color.pdf}}
}   

\newpage\setcounter{page}{1}
\pagenumbering{arabic}


\newcommand{\twopartdef}[4]
{
	\left\{
		\begin{array}{ll}
			#1 & \mbox{si } #2 \\
			#3 & \mbox{si } #4
		\end{array}
	\right.
}


\begin{document}


\section*{Gu\'ia Pr\'actica 4 -- Complejidad}

\begin{enumerate}
	\item Mostrar usando la definición de $\mathcal{O}()$ tomada de \footnote{Marzal y Gracia. Introducción al análisis de algoritmos.} \footnote{\href{https://arco.esi.uclm.es/public/mirror/Introducci\%C3\%B3n\%20al\%20an\%C3\%A1lisis\%20de\%20algoritmos.pdf}{https://arco.esi.uclm.es/public/mirror/Introducci\%C3\%B3n\%20al\%20an\%C3\%A1lisis\%20de\%20algoritmos.pdf}} : \\
	\vspace*{-7mm}
	\begin{center}
		$\mathcal{O}(f(n)) = $ \{ $t: \mathbb{N} \rightarrow \mathbb{R}^+$ $|$  $\exists c \in \mathbb{R}^+, 
		\exists n_0 \in \mathbb{N} : \forall n \geq n_0$  $\Rightarrow$   $t(n) \leq c f(n)$ \}.
	\end{center}
	
		\begin{enumerate}[a)]
			\item $50n-8 \in \mathcal{O}(n)$ 
			\item $5n^2-3n + 6 \in \mathcal{O}(n^2)$
			\item $2^n + 7n + 3 \in \mathcal{O}(2^n)$ 			
		\end{enumerate}
		
	\item Determinar $\mathcal{O}()$ y $\Omega()$ de las siguientes funciones:
		\begin{enumerate}[a)]
			\item $5n^2$ 
			\item $3n^2 + 2n$ 			
			\item $3^n + 2^n$ 
			\item $n + \log_2(n)$ 			
			\item $n^2 + \log_2(n)$ 			
		\end{enumerate}	
		
	\item Determinar $\Theta()$ de las siguientes funciones:
		\begin{enumerate}[a)]
			\item $n^2 + 2n$ 
			\item $t(n) = \twopartdef { 2n + 4 } {$ n  par $} {7n^2 + n} {$n impar$}$
			\item $t(n) = \twopartdef { 4n + 5 } {$ n  par $} {5n + n} {$n impar$}$
		\end{enumerate}	
		
	\item Reescribir los siguientes órdenes de complejidad de menor a mayor ($<$):\\
	$\mathcal{O}(n), \mathcal{O}(n\log n), \mathcal{O}(n!), \mathcal{O}(n^2), \mathcal{O}(5),
	\mathcal{O}(\log\log n), \mathcal{O}(3^n), \mathcal{O}(\sqrt{n}), \mathcal{O}(n^n),$ \\
	 $\mathcal{O}(n^3), \mathcal{O}(2^n), \mathcal{O}(\log n)$. 
	 
	 Etiquete cada $\mathcal{O}$ como
	 \textbf{sublineal}, \textbf{lineal} o \textbf{superlineal}. Dentro de éstas últimas puede
	 decir cuáles son \textbf{polinómicas} (cuadráticas, cúbicas, etc.) o \textbf{exponenciales}.
		
		
	\item Para cada uno de los fragmentos de código siguientes:
	\vspace{-3mm}
	\begin{itemize}
		\item Diga qué $ \mathcal{O}()$ poseen.
		\item Impleméntelos y dé el tiempo de ejecución para $\neq$ valores de $n$.
		\item Compare estimación teórica (primer ítem) con tiempos reales (segundo ítem).
	\end{itemize}
	
\newpage

\begin{minipage}[t]{7cm}
\null 
 \begin{algorithm}[H] %1
\dontprintsemicolon
$sum \gets 0$\;
\For{$i \gets 1$ \textbf{to} $n$} {
	$sum \gets sum + 1$
}
\caption{{ } }
 \end{algorithm}
\end{minipage}%
\begin{minipage}[t]{7cm} %2
\null
 \begin{algorithm}[H]
\dontprintsemicolon
$sum \gets 0$\;
\For{$i \gets 1$ \textbf{to} $n$} {
	\For{$j \gets 1$ \textbf{to} $n$} {
		$sum \gets sum + 1$
	}
}
\caption{{} } 
  \end{algorithm}
\end{minipage}


\vspace*{15mm}

\begin{minipage}[t]{7cm}
\null 
 \begin{algorithm}[H]
\dontprintsemicolon %3
$sum \gets 0$\;
\For{$i \gets 1$ \textbf{to} $n$} {
	\For{$j \gets 1$ \textbf{to} $n^2$} {
		$sum \gets sum + 1$
	}
}
\caption{{} } 
  \end{algorithm}
\end{minipage}%
\begin{minipage}[t]{7cm}
\null
 \begin{algorithm}[H]
\dontprintsemicolon%4
$sum \gets 0$\;
\For{$i \gets 1$ \textbf{to} $n$} {
	\For{$j \gets 1$ \textbf{to} $i$} {
		$sum \gets sum + 1$
	}
}
\caption{{} }
  \end{algorithm}
\end{minipage}


\vspace*{15mm}


\begin{minipage}[t]{7cm}
\null 
 \begin{algorithm}[H]
\dontprintsemicolon %5
$sum \gets 0$\;
\For{$i \gets 1$ \textbf{to} $n$} {
	\For{$j \gets 1$ \textbf{to} $i^2$} {
		\For{$k \gets 1$ \textbf{to} $j$} {
			$sum \gets sum + 1$
		}
	}
}
\caption{{} }
  \end{algorithm}
\end{minipage}%
\begin{minipage}[t]{7cm}
\null
 \begin{algorithm}[H]
\dontprintsemicolon %6
$i \gets 1$\;
\While{$i \leq n$} {
	$j \gets 1$\;
	\While {$j \leq n$} {
		$j \gets j*2$
	}
	$i \gets i + 1$
}
\caption{{ } }
\end{algorithm}
\end{minipage}


\vspace*{15mm}


\begin{minipage}[t]{7cm}
\null 
 \begin{algorithm}[H]
\dontprintsemicolon %7
$i \gets 1$\;
\While{$i \leq n$} {
	$j \gets 1$\;
	\While {$j \leq i$} {
		$j \gets j+1$
	}
	$i \gets i + 1$
}
\caption{{} }
  \end{algorithm}
\end{minipage}%
\begin{minipage}[t]{7cm}
\null
 \begin{algorithm}[H]
\dontprintsemicolon %8
$i \gets 1$\;
\While{$i \leq 10$} {
	$j \gets 1$\;
	\While {$j \leq 10$} {
		$j \gets j+1$
	}
	$i \gets i + 2$
}
\caption{{ } }
\end{algorithm}
\end{minipage}

\newpage
%%%%%

	\item ¿Cuál es el propósito de los siguientes algoritmos? Analizar el tiempo de
	ejecución en el peor caso y expresarlos en notación $\mathcal{O}$.
		
\begin{minipage}[t]{7cm}
\null 
 \begin{algorithm}[H]
\dontprintsemicolon %1
\KwIn{$a, n$: \textbf{integer}}
\KwOut{b}
$b$: \textbf{long integer}\;
$b \gets 1$\;
\For{$k \gets 0$ \textbf{to} $k < n$} {
$b \gets b*a$
}
\caption{{} }
  \end{algorithm}
\end{minipage}%
\begin{minipage}[t]{7cm}
\null
 \begin{algorithm}[H]
\dontprintsemicolon %2
\KwIn{$a, n$: \textbf{integers}}
\KwOut{b}
$b$: \textbf{long integer}\;
$c$: \textbf{long integer}\;
$b \gets 1$\;
$c \gets a$\;
\For{$k \gets n$ \textbf{to} $k > 0$} {
  \If{$k$ \textbf{mod} 2 == $0$} {
	$k \gets k / 2$\;
	$c \gets c*c$\;
  }
  \Else{
	$k \gets k -1$\;
	$b \gets b*c$\;
  }
}
\caption{{ } }
\end{algorithm}
\end{minipage}

\vspace{15mm}		
		
		
		\item A partir de una secuencia de números $A$, calcular otra $B$ tal que cada uno
		de sus elementos $B_i$ sea el promedio de todos los anteriores en la secuencia
		original, incluído el $i-ésimo$.\\
		Ejemplo:\\
		$A = [1, 3, 5, 4, 2, 1]$\\
		$B = [1, 2, 3, 3.25, 3, 2.66]$
		
		
		\item Un arreglo $A$ contiene $n-1$ enteros \textbf{únicos} $\in [0, n-1]$; es decir, 
		hay un número de este intervalo que no está en $A$. Escriba un algoritmo para encontrar
		ese número faltante en $\mathcal{O}(n)$. Sólo es posible utilizar $\mathcal{O}(1)$ espacio
		adicional además del arreglo $A$ mismo.\\
		Ejemplo:\\
		$N = 5$. El arreglo puede contener elementos (no repetidos) $\in [0, 4]$.\\
		Si el vector $=[1, 4, 2, 0]$, entonces el elemento que falta es \textbf{3}.
		

		\item Suponga que cada renglón de un arreglo $A_{n,n}$ posee valores $0$ y $1$ de tal
		manera que en cualquier renglón de $A$ todos los $1$'s van antes que los $0$'s. Describa
		un método de $\mathcal{O}(n)$ para determinar el renglón de $A$ que contenga la
		mayor cantidad de 1's.\\\\
		Ejemplo:
		\begin{flushleft}
  \begin{tabular}{|| c | c | c | c | c ||}
    \hline\hline
    1 & 1 & 0 & 0 & 0 \\ \hline
    1 & 1 & 1 & 0 & 0 \\ \hline
    \cellcolor{green!17}\textbf{1} & \cellcolor{green!17}\textbf{1} & \cellcolor{green!17}\textbf{1} & \cellcolor{green!17}\textbf{1} & \cellcolor{green!17}\textbf{0} \\ \hline
    1 & 0 & 0 & 0 & 0 \\ \hline
    1 & 1 & 1 & 0 & 0 \\ \hline
    \hline
  \end{tabular}
\end{flushleft}


	\item Para cada solución que tenga de los problemas de las guías anteriores, comente, si es posible, el $\mathcal{O}()$ de
	su solución. Si no es posible (por ejemplo talla indeterminada del problema, o caminos de ejecución no determinísticos) explique
	por qué. Para los problemas que no tenga resuelto, trate de proponer un algoritmo (al menos en forma coloquial) y vea
	si puede determinar el $\mathcal{O}()$ de su propuesta.
	
	\item Realice algoritmos (no es necesario implementarlos, suficiente con pseudocódigo claro) para:
		\begin{enumerate}[a)] 
			\item \textbf{Sumar} $N \in \mathbb{N}$ y  $M \in \mathbb{N}$ tales que $n$ y $m$ representan la cantidad de dígitos cada uno. Piense que el número
			de dígitos puede ser muy grande.
			\item \textbf{Multiplicar} $N \in \mathbb{N}$ y  $M \in \mathbb{N}$ tales que $n$ y $m$ representan la cantidad de dígitos cada uno. Piense que el número
			de dígitos puede ser muy grande.
			\item \textbf{Multiplicar} las \textbf{matrices} $A_{m,n}$ y $B_{n,q}$.
			
			\item[] \underline{\textbf{NOTA}}: En todos los casos exprese el $\mathcal{O}$() de su algoritmo.
		\end{enumerate}
	
	\item Resolver en lenguage C++ y obtener \textbf{ACCEPTED} en los jueces correspondientes para los siguientes problemas:	
	
	\begin{enumerate}[a)]
		\item Integer Inquiry  -- \href{http://uva.onlinejudge.org/external/4/424.html}{http://uva.onlinejudge.org/external/4/424.html}
		\item 500!  -- \href{http://uva.onlinejudge.org/external/6/623.html}{http://uva.onlinejudge.org/external/6/623.html}
		\item Overflow -- \href{http://uva.onlinejudge.org/external/4/465.html}{http://uva.onlinejudge.org/external/4/465.html} 
		
		\item[] \underline{\textbf{NOTA}}: En todos los casos exprese el $\mathcal{O}$() de su algoritmo.
	\end{enumerate}


\end{enumerate}

\end{document}
